
\begin{proposition}[Proved exclusion of Enoki components]\label{cdp:st:C.2.INT.2}

Let $f:Y\to B$ be a proper connected-fibre map from a
smooth compact complex threefold with $H^2(Y,\Q)=0$
and general fibre a complex two-torus. No component
of a reducible fibre has normalization resolution
with Enoki minimal model.
\end{proposition}
\begin{proof}
Write the full fibre as $F=\sum_i m_iS_i$. Its
normalization resolutions $Z_i$ are non-K\"ahler by
\ref{cdp:st:C.3.4}. A neighbouring component gives on each $Z_i$
a nonzero effective divisor $E_{ij}=h_i^*S_j$ of
zero real Chern class. On a class VII minimal model
the precise classification in \ref{cdp:st:C.2.INT.1} makes that
model Hopf or Enoki. The other non-K\"ahler minimal
models are elliptic surfaces, including both types
of Kodaira surfaces. \ref{cdp:st:C.2.INT.1} rules out an interface
between either such type and an Enoki component.
Thus connectedness of the fibre shows that, if an
Enoki component occurs, every component has Enoki
minimal model.

Let $D_i^{\rm self}\ge0$ be the effective resolution
conductor for $S_i$ from \ref{cdp:st:A.5.3}, and set
$D_i=D_i^{\rm self}+\sum_{j\ne i}E_{ij}$. This is
also the resolution conductor for the total reduced
support. Adjunction specifies the exact comparisons
\[
 K_{Z_i}+D_i=h_i^*(K_Y+\sum_j S_j),\qquad
 K_{Z_i}+D_i^{\rm self}=h_i^*(K_Y+S_i).
 \tag{C.INT.5}\label{cdp:eq:C.INT.5}
\]
The effective numerically anticanonical divisor on
an Enoki minimal model makes it parabolic Inoue;
this is the Enoki case in
\cite[Lemma~2.22 and the following paragraph]{Dloussky2006}.
Write $E_i^0$ for its elliptic curve and $C_i$ for
its reduced rational cycle. They are disjoint and
\[
 K_{Z_i^0}+E_i^0+C_i=0.
 \tag{C.INT.6}\label{cdp:eq:C.INT.6}
\]
The pushforward of $D_i$ has the form
$E_i^0+k_i' C_i$, with $k_i'\ge0$, by intersection
with $E_i^0$ and the kernel of the cycle intersection
matrix. Each nonzero adjacent divisor has the exact
form $E_{ij}=n_{ij}\sigma_i^*C_i$, $n_{ij}\ge1$.
In particular $k_i'\ge\sum_{j\ne i}n_{ij}\ge1$.

The canonical evaluation used in \ref{cdp:st:A.4.4} gives, on a
neighbourhood of this fibre,
$K_Y=\OO_Y(\sum_i k_iS_i)$ for integers $k_i\ge0$;
the line from the base is trivial on the chosen disc.
Put $r_i=(k_i+1)/m_i$. The principal full-fibre
relation yields the following equality in
$\Pic(Z_i)\otimes\Q$:
\[
 K_{Z_i}+D_i
   =\sum_{j\ne i}m_j(r_j-r_i)E_{ij}.
 \tag{C.INT.7}\label{cdp:eq:C.INT.7}
\]
Indeed subtract $r_i h_i^*F=0$ from the first
formula \eqref{cdp:eq:C.INT.5}. Choose $i$ for which $r_i$ is
maximal. The negative of the right side of
\eqref{cdp:eq:C.INT.7} is an effective rational divisor $A_i$.
Thus $K_{Z_i}+D_i+A_i$ is torsion after clearing
denominators. Push to its minimal model and use
\eqref{cdp:eq:C.INT.6}. Taking Gauduchon degree gives
\[
 k_i'+\sum_{j\ne i}m_j(r_i-r_j)n_{ij}=1,
 \tag{C.INT.8}\label{cdp:eq:C.INT.8}
\]
since $\deg\OO(C_i)>0$. All terms after $k_i'$
are nonnegative, and $k_i'\ge\sum n_{ij}\ge1$.
Equality forces $k_i'=1$, exactly one neighbour,
$n_{ij}=1$, and $r_j=r_i$. That neighbour is again
a maximum. Propagating along the connected
intersection graph proves that there are exactly
two components $S_1,S_2$, all ratios are equal,
and on both minimal models
\[
 (D_i^{\rm self})_0=E_i^0,
 \qquad (E_{ij})_0=C_i.
 \tag{C.INT.9}\label{cdp:eq:C.INT.9}
\]

We retain the geometric information in the actual
blowups and normalization maps. Every topologically
torsion line descends through a point blowup, by the
integral Picard blowup formula. Applying the
canonical divisor formula to the second equality
\eqref{cdp:eq:C.INT.5} therefore gives at every blowup
\[
 (D_i^{\rm self})_{\rm new}
   =\sigma^*(D_i^{\rm self})_{\rm old}-E_{\rm exc}.
 \tag{C.INT.10}\label{cdp:eq:C.INT.10}
\]
Starting with the smooth multiplicity-one elliptic
curve $E_i^0$, effectivity permits centres only on
its current strict transform, and the new
exceptional coefficient is zero. Hence
$D_i^{\rm self}$ is that elliptic strict transform,
and all these centres are disjoint from $C_i$.
The external cycle $\sigma_i^*C_i=C_i$ is therefore
unchanged, and disjoint from the self conductor.

It follows that $S_i$ is normal near the external
intersection. Here is the needed resolution
comparison at possible exceptional points. Any
resolution-exceptional curve outside $D_i$ is
disjoint from it: adjunction and minimality give
$K_{Z_i}\cdot E\ge0$, while $D_i\cdot E\ge0$,
so $(K_{Z_i}+D_i)\cdot E=0$ forces $D_i\cdot E=0$.
Connected resolution fibres cannot join the
elliptic self conductor to the external rational
cycle. The conductor on the normal surface is a
reflexive ideal, by the duality calculation in
\ref{cdp:st:A.5.3}, so it has no isolated additional support
on that cycle. Thus normalization is an isomorphism
near its image. Resolution over this image is
crepant by the second formula \eqref{cdp:eq:C.INT.5}, since
$D_i^{\rm self}$ is disjoint from it. Its
exceptional curves are consequently $(-2)$
curves in $C_i$. A whole cycle cannot be contracted,
because its intersection matrix is not negative
definite. The connected fibres which are
contracted are proper chains in that cycle.

The intersection scheme $D=S_1\cap S_2$ is a
complete-intersection curve in $Y$. Its pullback
has reduced divisor $C_i$, so it is generically
reduced; its Cohen--Macaulay property then makes
it reduced everywhere. It is connected, since it
is the image of $C_i$. The equality of all $r_i$
means $K_Y+S_1+S_2=rF$ in $\Pic(Y)\otimes\Q$
on the fibre neighbourhood. Adjunction gives
$\omega_D$ torsion, hence $p_a(D)=1$.
Topologically the map $C_i\to D$ contracts just
the proper chains described above and is an
isomorphism elsewhere. It preserves the graph
circle and all its remaining components are
rational. Thus its normalized components are
$\PP^1$, its graph has exactly one cycle, and
$\sum\delta_x=\sum(r_x-1)$. Each contracted
chain has only two ends; every singular point
therefore has two branches and $\delta=1$.
The colength-one ring comparison (C.VII.6)
proves that these are ordinary nodes. At least
one node occurs.

At such a node the complete-intersection curve
has embedding dimension two. At least one of
the two surface equations has a nonzero linear
part; take it to be $z$. On $z=0$ the other
equation has an ordinary node. The holomorphic
Morse lemma with parameter $z$ therefore writes
it as $xy+h(z)$. Normality of that surface near
$D$ excludes $h\equiv0$, and a change of the
$z$ coordinate writes $h(z)=z^n$, $n\ge1$.
Consequently the original fibration has local
equation
\[
 f=u(x,y,z)\,z^a(xy+z^n)^b,
 \qquad a,b,n\ge1,
 \tag{C.INT.11}\label{cdp:eq:C.INT.11}
\]
where $u$ is a unit and $a,b$ are the two full
fibre multiplicities. The unit can be removed
by an actual biholomorphic change of variables.
Put $d=2(a+nb)$ and choose a local logarithm of
$u$. For $f_0=z^a(xy+z^n)^b$, weighted homogeneity
gives
\[
 \Phi(x,y,z)=(x u^{n/d},y u^{n/d},z u^{2/d}),
 \qquad f_0\circ\Phi=u f_0.
 \tag{C.INT.12}\label{cdp:eq:C.INT.12}
\]
The derivative of $\Phi$ at the origin is
invertible, proving the required equivalence
of the actual local fibres.

For $t\ne0$ small, take a simply connected
sector in the punctured $z$-plane and one
branch of $(t z^{-a})^{1/b}$. An open part
of the actual smooth fibre is
\[
 xy=g(z):=(t z^{-a})^{1/b}-z^n.
 \tag{C.INT.13}\label{cdp:eq:C.INT.13}
\]
Consecutive zeros on this logarithmic branch
have angular separation $2\pi b/(a+nb)<2\pi$;
two such zeros fit in one sector. They are
simple, since at either one
$g'(z)=-(a/b+n)z^{n-1}\ne0$.
Join them by a smooth embedded path with no
other zero and take the circle $|x|=|y|$
over its interior. Explicitly, for a continuous
choice of $\arg g$ along that interior, it is
\[
 x=|g|^{1/2}e^{i(\arg g/2+\theta)},\qquad
 y=|g|^{1/2}e^{i(\arg g/2-\theta)}.
 \tag{C.INT.14}\label{cdp:eq:C.INT.14}
\]
The circle collapses to a smooth disc at each
simple zero, giving an embedded sphere.
The two $x,y$ terms of the standard K\"ahler
form cancel on this sphere, and the $z$ term
vanishes on the one-dimensional path. The
sphere is Lagrangian; its normal bundle is
$JTS^2$ with the opposite Euler orientation,
so its self-intersection is $-2$.
The whole sphere lies in the local chart:
$|z|=O(|t|^{1/(a+nb)})$ and
$|x|,|y|=O(|t|^{n/(2(a+nb))})$.
Its image in the compact smooth torus fibre
would be null-homotopic, because a torus has
contractible universal cover. Its integral
homology class would therefore be zero and
its self-intersection zero, contradicting
$-2$. This contradiction excludes the Enoki
alternative.
\end{proof}


